%!TEX program = lualatex
\documentclass[11pt, a4paper]{scrartcl}
% Packages
\usepackage[margin=1.25in]{geometry}
\usepackage{index}
\usepackage{amsbsy} % Bold math symbols
\makeindex
\usepackage{tcolorbox}
\usepackage{varwidth}
\usepackage{amsmath, amssymb}
\usepackage{esint}
\usepackage{titlesec}
\usepackage{xcolor}
\usepackage{titling}
\usepackage[linktocpage]{hyperref}
\usepackage{pgfplots}
\usepackage{multicol}
\setlength{\columnsep}{2em}
\usepackage{caption}
\usepackage{amsthm}
\usepackage{import}
\usepackage{cancel}
\usepackage{caption}
\usepackage{nicematrix}
\usepackage{mathtools}
\usepackage{enumerate}
\usepackage{graphicx}
\usepackage{makecell}
\usepackage{tikz}
\usepackage{circuitikz}
\usepackage{pgf-spectra}
\usepackage[version=4]{mhchem}
\usetikzlibrary{arrows.meta,calc,positioning, angles,quotes,patterns.meta,intersections}
\usepackage[italian]{babel}
% Font
%\usepackage{fontspec}
%\setmainfont{Libertinus Serif}

%\usepackage{unicode-math}
%\setmathfont{Libertinus Math}


% To reset footnote numbering each page
\usepackage[perpage]{footmisc}
\usepackage{setspace}
\setstretch{1.2}

% Titles 
\title{Note di Fluidodinamica}
\author{}
\date{}


% svolgimento
\newenvironment{svolgimento}{\renewcommand\qedsymbol{$\blacksquare$}\begin{proof}[Svolgimento]}{\end{proof}}

\newtheoremstyle{style}% name of the style to be used
{5pt}% measure of space to leave above the theorem. E.g.: 3pt
{5pt}% measure of space to leave below the theorem. E.g.: 3pt
{\normalfont}% name of font to use in the body of the theorem
%{15pt}% measure of space to indent
{0pt}% measure of space to indent
{\noindent\bfseries}% name of head font
{}% punctuation between head and body
{ }% space after theorem head; " " = normal interword space
{\thmname{#1}\thmnumber{ #2}{\thmnote{ (#3)}.\ }}


\theoremstyle{style}
\newtheorem{esempio}{Esempio}[section]
\newtheorem{definizione}{Definizione}[section]
\newtheorem{prop}{Proposizione}[section]
\newtheorem{teorema}{Teorema}[section]
\newtheorem{lemma}{Lemma}[teorema]
\newtheorem{corollario}{Corollario}[teorema]
\newtheorem{osservazione}{Osservazione}[section]
\newtheorem{assunzione}{Assunzione}[section]
\newtheorem{notazione}{Notazione}[section]

\newtheoremstyle{style1}% name of the style to be used
{5pt}% measure of space to leave above the theorem. E.g.: 3pt
{5pt}% measure of space to leave below the theorem. E.g.: 3pt
{\normalfont}% name of font to use in the body of the theorem
%{15pt}% measure of space to indent
{0pt}% measure of space to indent
{\noindent\bfseries\color{red}}% name of head font
{}% punctuation between head and body
{ }% space after theorem head; " " = normal interword space
{\thmname{#1}\thmnumber{ #2}{\thmnote{ (#3)}.\ }}

\theoremstyle{style1}
\newtheorem{esercizio}{Esercizio}[section]


%% Generic box
\newtcolorbox{eqbox}[1][]
{
colback=gray!10,
arc=0pt,
boxrule=0pt,
title=#1
}

 \newenvironment{boxenv}[1][]{
    \begin{eqbox}[#1]
    }{
   \end{eqbox}
}



%%%%%%%%%% Medie con integrali multipli
\def\Yint#1{\mathchoice
    {\YYint\displaystyle\textstyle{#1}}%
    {\YYint\textstyle\scriptstyle{#1}}%
    {\YYint\scriptstyle\scriptscriptstyle{#1}}%
    {\YYint\scriptscriptstyle\scriptscriptstyle{#1}}%
      \!\iint}
\def\YYint#1#2#3{{\setbox0=\hbox{$#1{#2#3}{\iint}$}
    \vcenter{\hbox{$#2#3$}}\kern-.51\wd0}}
\def\longdash{{-}\mkern-3.5mu{-}} 
   % consider using "\mkern-7.5mu" if esint package is loaded
\def\tiltlongdash{\rotatebox[origin=c]{15}{$\longdash$}}
\def\fiint{\Yint\tiltlongdash}

\def\Zint#1{\mathchoice
    {\YYint\displaystyle\textstyle{#1}}%
    {\YYint\textstyle\scriptstyle{#1}}%
    {\YYint\scriptstyle\scriptscriptstyle{#1}}%
    {\YYint\scriptscriptstyle\scriptscriptstyle{#1}}%
      \!\iiint}
      \def\tilongdash{\mkern6mu{-}\mkern-4mu{-}\mkern-5mu{-}} 
   % consider using "\mkern-7.5mu" if esint package is loaded
\def\titiltlongdash{\rotatebox[origin=c]{15}{$\tilongdash$}}
\def\fiiint{\Zint\titiltlongdash}

%Captions
\captionsetup[figure]{font=footnotesize,labelfont=footnotesize}
\captionsetup[table]{font=footnotesize,labelfont=footnotesize}
%Titlesec
\titleformat{\section}
{\fontsize{15}{20}\sffamily\scshape}
{\normalfont\color{gray}{\fontsize{20}{20}\selectfont\thesection}}
{0.7em}
{}
\hypersetup{colorlinks,breaklinks, linkcolor=[RGB]{74, 122, 164}}
\definecolor{asdf}{HTML}{4a7aa4}
% Personalizza la formattazione della subsection
\titleformat{\subsection}[block]{\fontsize{13}{20}\bfseries}{\normalfont\thesubsection}{.5em}{}


% Personalizza la formattazione della subsubsection
\titleformat{\subsubsection}[block]{\fontsize{12}{20}\bfseries}{\normalfont\thesubsubsection}{.5em}{}

% Maketitle customization
\renewcommand{\maketitle}{
\begin{center}
{\sffamily
{\fontsize{20}{20}\selectfont\MakeUppercase\thetitle}}

\vspace{0.2in}

{\large\scshape\sffamily\theauthor}
\end{center}
}


% Differenziale 
\newcommand{\dd}{\mathop{}\!\mathrm{d}}




%Evaluate symbol
\DeclareMathOperator{\di}{d\!}
\newcommand*\Eval[3]{\left.#1\right\rvert_{#2}^{#3}}

%%%%%%% Numero delle equazioni in formato a.b
\numberwithin{equation}{subsection}
%%%%%

%%%%%%%%%% Personalizzazione numeri lista
\renewcommand{\theenumi}{(\arabic{enumi})}

%%%% Table of contents

\usepackage[titles]{tocloft}

\renewcommand{\cftdot}{}
\usepackage{titletoc}
%\setcounter{tocdepth}{2}

%%%%%%%%%%%%%%%% Toc style

% Personalizzazione scritta indice


%%%% Lambda bar
\newcommand{\lbar}{{\mkern0.75mu\mathchar '26\mkern -9.75mu\lambda}}


\begin{document}
\maketitle
\newpage
\tableofcontents 
\newpage
\section{Fondamenti teorici}

\subsection{Basi del modello}
La descrizione del moto di un fluido si basa sulla descrizione del moto di alcune sue particelle; in particolare, si studia il comportamento di un insieme di queste sue particelle, detto \textit{elemento fluido}.
Si assegna, a questo elemento fluido, una certa massa $m$ e lo si studia come fosse un punto materiale:
\[
\mathbf{F}  = m\mathbf{a}  \hspace{3cm} \begin{cases}
	\displaystyle \frac{d \mathbf{x} }{d t} = \mathbf{v} \big(\mathbf{x} (t) , t\big)\\
	\mathbf{x} (0) = \mathbf{x} _0
\end{cases} 
\] 
Questa trattazione presuppone che non vi sia necessit\`a di specificare altre informazioni sull'elemento fluido che si sta studiando.
Di fatto, questo elemento fluido avr\`a un suo volume, cio\`e una sua forma, e, durante l'evoluzione temporale, non \`e assicurato che questo rimanga uguale.
Per il momennto, si assume, per\`o, che non vi sia necessit\`a di considerare questi aspetti.

La legge di conservazione della quantit\`a di moto che deriva direttamente da $0=\mathbf{F}  = m\mathbf{a} $ in assenza di forze \`e valida purch\'e non vi sia un sostanziale scambio di particelle con altri elementi fluidi.
Di fatto, col passare del tempo, le particelle che compongono l'elemento fluido in esame tenderanno ad essere \textit{riciclate}; quindi, \`e importante tenere in considerazione che, passato sufficientemente tempo, dipendente dal volume dell'elemento fluido, la conservazione del momento potrebbe non essere pi\`u soddisfatta per via del fatto che le particelle che compongono l'elemento fluido cambiano.

Oltretutto, la validit\`a della descrizione fluidodinamica del sistema in esame sar\`a strettamente legata, anche per quanto appena detto, a una lunghezza e un tempo caratteristici, rispettivamente $\ell$, $ \tau $: prendendo un elemento fluido con $\ell $ \textit{troppo piccolo}, $\tau $ sar\`a conseguentemente troppo piccolo, rendendo impossibile una descrizione del sistema.
Il tempo caratteristico $\tau $ rappresenta, in un certo senso, il tempo di vita dell'elemento fluido, cio\`e quel tempo oltre il quale la trattazione fluidodinamica perde di validit\`a.
\subsection{Stima della validit\`a del modello}

Si vuole dare una stima per capire quando una trattazione fluidodinamica \`e possibile e valida.

Si considerano $n$, $\sigma $, $\ell _c$, $\tau _c$, rispettivamente, densit\`a numerica di particelle, sezione d'urto relativa agli urti tra le particelle dell'elemento fluido, cammino libero medio e tempo medio tra un urto e l'altro.
\begin{osservazione}
Questa trattazione \`e intrinsecamente fuori dal dominio della fluidodinamica, in quanto si stanno considerando grandezze microscopiche della materia.
\end{osservazione}
\noindent Per definizione, si ha $n\sigma \ell _c = 1 $ e il tempo medio di percorrenza \`e ottenuto come $\ell _c = v_\text{th} \tau _c$, dove $v_\text{th}$ \`e la velocit\`a termica delle particelle.
La \textit{legge di Fick}\footnote{Fondamentalmente, questa afferma che la variazione temporale di particelle per unit\`a di volume \`e legata alla distribuzione delle particelle tramite il coefficiente di diffusione: concentrando un certo numero di particelle in un volume, si osserva la tendenza di queste ultime a distribuirsi omogeneamente nei dintorni, senza generare squilibri; l'andamento temporale della densit\`a \`e legato a quanto intenso \`e lo squilibrio.} afferma che
\[ 
\partial _t \rho = D \nabla ^2 \rho 
\] 
dove $\rho $ densit\`a (numerica o massiva a seconda che le particelle in questione abbiano una massa ben definita o meno) e $D$ coefficiente di diffusione.
Si utilizza questa legge per stimare il tempo necessario perch\'e l'elemento fluido sostituisca in modo sensibile le sue particelle.
Essendo interessati ad una stima, si pu\`o scrivere la legge di Fick in termini di rapporto incrementale: definendo un tempo caratteristico $T$ rispetto a cui la densit\`a di particelle nell'elemento fluido \`e nulla (cio\`e le particelle che componevano l'elemento fluido originario sono state tutte sostituite) e definendo una lunghezza caratteristica $L $ associata al tempo $T$, si pu\`o scrivere che
\[
\frac{\partial \rho }{\partial t} \sim \frac{\rho }{T}\, \, , \ \ \nabla ^2 \rho \sim \frac{\rho }{L^2}\implies \frac{1}{T} = \frac{D}{L ^2}
\] 
Ne segue che, affinch\'e la trattazione fluidodinamica sia valida, deve valere 
\begin{equation}
	\tau \ll \frac{L^2}{D}
\end{equation}
Per un gas in condizioni standard, si pu\`o esprimere il coefficiente di diffusione tramite i parametri microscopici riportati sopra:
\[
D = \frac{\ell_c ^2}{\tau _c} = \ell _c v_{\text{th}} \approx 1.4 \times 10^{-5} \text{ (MKS)}
\] 





























\end{document}
